% ============================================================
%  UniVerify — IEEE Conference Paper  (4–5 pages)
% ============================================================
\documentclass[conference]{IEEEtran}
\IEEEoverridecommandlockouts

\usepackage[T1]{fontenc}
\usepackage[utf8]{inputenc}
\usepackage{amsmath,amssymb}
\usepackage{booktabs}
\usepackage{xcolor}
\usepackage{pgfplots}
\usepackage{tikz}
\usepackage{microtype}
\usepackage{cite}
\usepackage{url}
\usepackage{float}
\usepackage{multirow}
\usepackage{subfig}
\usepackage{xfp}
\usepackage{array}

\pgfplotsset{compat=1.18}
\usetikzlibrary{
  positioning, arrows.meta, shapes.rounded,
  fit, backgrounds, calc, decorations.pathreplacing
}

% ── Palette ───────────────────────────────────────────────────────────────
\definecolor{clrVerified}{HTML}{4CAF50}
\definecolor{clrProof}   {HTML}{FF9800}
\definecolor{clrComputed}{HTML}{2196F3}
\definecolor{clrRoot}    {HTML}{1B5E20}
\definecolor{clrGrey}    {HTML}{CFD8DC}
\definecolor{clrOnChain} {HTML}{E3F2FD}
\definecolor{clrOffChain}{HTML}{FFF8E1}
\definecolor{clrCard}    {HTML}{ECEFF1}

% ── Cost macro ────────────────────────────────────────────────────────────
\newcommand{\EthUsd}{2500}
\newcommand{\CostUsd}[2]{\fpeval{round(#1*(#2*1e-9)*(\EthUsd),4)}}

% ── Compact table rows ────────────────────────────────────────────────────
\renewcommand{\arraystretch}{1.05}

% =============================================================
\begin{document}

\title{Gas-Efficient Diploma Registry on Ethereum:\\
A Three-Phase Cost Analysis}

\author{
  \IEEEauthorblockN{Nazar Shcherbyna}
  \IEEEauthorblockA{
    Faculty of Electronics and Information Technology\\
    Warsaw University of Technology, Warsaw, Poland\\
    \texttt{nazar.shcherbyna@student.pw.edu.pl}
  }
}

\maketitle

% =============================================================
\begin{abstract}
Issuing tamper-proof academic credentials on Ethereum is attractive,
but per-transaction gas costs can make large-scale deployment
prohibitively expensive.
We present UniVerify, a Solidity diploma registry studied across
three sequential optimisation phases.
First, seven code-level changes to a per-diploma model reduce issuance
cost by 14\% and deployment cost by 20\%.
Second, replacing per-diploma storage with a Merkle batch architecture
cuts per-diploma issuance cost by \textbf{98.9\%} at $N=100$
through amortisation of a single on-chain root write.
Third, four EVM hot-path optimisations yield savings captured by the
closed-form model $\Delta(d)=2^d\!\cdot\!42+96+d\!\cdot\!199$\,gas.
Our results show that architectural decisions outweigh code-level
tuning by roughly two orders of magnitude.
\end{abstract}

\begin{IEEEkeywords}
blockchain, Ethereum, Solidity, academic credentials,
Merkle tree, gas optimisation, smart contracts
\end{IEEEkeywords}

% =============================================================
\section{Introduction}

Academic credential fraud affects millions of graduates and
employers worldwide~\cite{grech2017blockchain}.
Blockchain registries offer a compelling remedy: once a diploma
fingerprint is stored on-chain, any party can verify it
without trusting a central authority.

The practical barrier is cost.
Every state-changing operation on Ethereum consumes \emph{gas},
priced in ETH at market rates.
Naively storing one hash per diploma can cost \$5--10 at moderate
gas prices — prohibitive for universities issuing thousands
of degrees annually.

Prior work focuses on the data model of blockchain credentials
(W3C DIDs, NFTs, IPFS anchoring) but rarely quantifies
the comparative economics of contract storage layouts
\cite{alketbi2018blockchain,lizcano2020decentralised}.
This paper fills that gap.

\medskip\noindent\textbf{Contributions:}
\begin{itemize}
  \item Systematic measurement of seven code-level optimisations
        on a per-diploma contract (§\ref{sec:phase1}).
  \item Quantitative amortisation model for a Merkle batch
        architecture (§\ref{sec:phase2}).
  \item Closed-form, experimentally validated savings model
        for EVM hot-path optimisations (§\ref{sec:phase3}).
\end{itemize}

% =============================================================
\section{UniVerify: System Design}
\label{sec:design}

\subsection{Two Storage Models}

\textbf{Per-diploma model.}
Each diploma is registered individually:
\texttt{issue(bytes32 docHash)} writes to a
\texttt{mapping(bytes32\,{\small$\Rightarrow$}\,uint256)}.
Revocation flips a bit in the same slot.

\textbf{Merkle batch model.}
The issuer builds a Merkle tree over the entire cohort \emph{off-chain}
(zero gas), then submits only the tree root via
\texttt{issueBatchRoot(uint64 batchId, bytes32 root)} —
one transaction regardless of cohort size $N$.
Membership verification is \texttt{statusWithProof(docHash, issuer, batchId, proof[])},
a \texttt{view} function executed as a free \texttt{eth\_call}.

\subsection{Leaf Construction}

Each leaf binds a diploma to its deployment context to prevent
cross-chain replay:
\begin{equation}
\ell_i = \mathrm{keccak256}(
  \texttt{this} \,\|\,
  \texttt{chainid} \,\|\,
  \texttt{issuer} \,\|\,
  \texttt{batchId} \,\|\,
  \texttt{docHash}_i).
\label{eq:leaf}
\end{equation}
Internal nodes use sorted pairs:
$H(a,b)=\mathrm{keccak256}(\min(a,b)\|\max(a,b))$,
which prevents second-preimage attacks without extra bookkeeping.

\subsection{Architecture and Merkle Tree}

Fig.~\ref{fig:overview} shows the complete workflow and the
structure of a 4-diploma Merkle batch with proof path highlighted.

% ── FIGURE 1: Architecture + Merkle tree (full-width) ───────
\begin{figure*}[t]
\centering
% ── Left panel: architecture ────────────────────────────────
\subfloat[System workflow. Tree construction is entirely off-chain;
only the root reaches the EVM.]{%
\begin{tikzpicture}[
  card/.style={
    draw=gray!60, fill=clrCard, rounded corners=6pt,
    minimum width=2.7cm, minimum height=0.65cm,
    font=\small, align=center, thick
  },
  zone/.style={
    draw=gray!50, rounded corners=8pt, dashed, thick,
    inner xsep=8pt, inner ysep=6pt
  },
  arr/.style={-{Stealth[length=6pt,width=5pt]}, thick},
  note/.style={font=\scriptsize, text=gray!80!black},
]
% ── Issuer zone (off-chain) ──────────────────────────────────
\node[card, fill=clrOffChain] (d1) at (0, 3.2) {Diploma 1 … $N$};
\node[card, fill=orange!25]   (bt) at (0, 2.1) {Build Merkle tree\\{\footnotesize (off-chain, \textbf{0 gas})}};
\node[card, fill=orange!40]   (rt) at (0, 1.0) {Root $r$};
\draw[arr] (d1)--(bt); \draw[arr] (bt)--(rt);

\begin{scope}[on background layer]
  \node[zone, fit=(d1)(bt)(rt),
    label={[font=\small\bfseries, yshift=2pt]above:Issuer (University)}] {};
\end{scope}

% ── Smart contract (on-chain) ────────────────────────────────
\node[card, fill=clrOnChain, minimum width=3.1cm, minimum height=0.9cm]
  (sc) at (3.4, 2.1)
  {\textbf{DiplomaRegistry}\\{\footnotesize\texttt{batchRoots[iss][id]=}$r$}};

\begin{scope}[on background layer]
  \node[zone, fill=clrOnChain!30, fit=(sc),
    label={[font=\small\bfseries]above:Ethereum (on-chain)}] {};
\end{scope}

% ── Verifier zone (off-chain) ────────────────────────────────
\node[card, fill=clrOffChain]  (denv) at (6.8, 3.2) {Diploma + proof};
\node[card, fill=green!20]     (call) at (6.8, 2.1) {\texttt{eth\_call}\\{\footnotesize\textbf{no gas charged}}};
\node[card, fill=clrVerified, text=white] (res) at (6.8, 1.0) {Valid / Revoked\\/ Unknown};
\draw[arr] (denv)--(call); \draw[arr] (call)--(res);

\begin{scope}[on background layer]
  \node[zone, fit=(denv)(call)(res),
    label={[font=\small\bfseries]above:Verifier (Employer)}] {};
\end{scope}

% ── Cross arrows ─────────────────────────────────────────────
\draw[arr, blue!70!black, line width=1.4pt]
  (rt) -- node[above, note]{\texttt{issueBatchRoot()}}
          node[below, note]{$\approx$51k gas} (sc);
\draw[arr, green!50!black, line width=1.4pt]
  (call) -- node[above, note]{\texttt{statusWithProof()}}
             node[below, note]{view, free} (sc);
\end{tikzpicture}%
\label{fig:arch}%
}
\hfill
% ── Right panel: Merkle tree ─────────────────────────────────
\subfloat[Merkle tree for 4 diplomas.
  \textcolor{clrVerified!80!black}{\textbf{Green}}: diploma under verification.
  \textcolor{clrProof!80!black}{\textbf{Orange}}: proof elements supplied by the holder.
  \textcolor{clrComputed!70!black}{\textbf{Blue}}: node recomputed by verifier.
  Only the root is stored on-chain.]{%
\begin{tikzpicture}[
  nd/.style={
    draw=gray!70, rounded corners=4pt, thick,
    minimum width=1.55cm, minimum height=0.62cm,
    font=\small, align=center
  },
  arr/.style={-{Stealth[length=5pt]}, thick, gray!60},
  prfarr/.style={-{Stealth[length=5pt]}, very thick},
]
% ── Leaves ───────────────────────────────────────────────────
\node[nd, fill=clrVerified, text=white]    (L0) at (0.00, 0) {$\ell_0$};
\node[nd, fill=clrProof]                   (L1) at (1.80, 0) {$\ell_1$};
\node[nd, fill=clrGrey]                    (L2) at (3.60, 0) {$\ell_2$};
\node[nd, fill=clrGrey]                    (L3) at (5.40, 0) {$\ell_3$};

\node[font=\scriptsize, align=center, text=clrVerified!70!black]
  at (0.00,-0.65) {verified\\diploma};
\node[font=\scriptsize, text=clrProof!80!black] at (1.80,-0.65) {proof{[0]}};
\node[font=\scriptsize, text=gray]              at (3.60,-0.65) {diploma 2};
\node[font=\scriptsize, text=gray]              at (5.40,-0.65) {diploma 3};

% ── Internal nodes ───────────────────────────────────────────
\node[nd, fill=clrComputed, text=white]    (H01) at (0.90, 1.70) {$H_{01}$};
\node[nd, fill=clrProof]                   (H23) at (4.50, 1.70) {$H_{23}$};

\node[font=\scriptsize, text=clrComputed!80!black] at (0.90, 1.08) {computed};
\node[font=\scriptsize, text=clrProof!80!black]    at (4.50, 1.08) {proof{[1]}};

% ── Root ─────────────────────────────────────────────────────
\node[nd, fill=clrRoot, text=white, minimum width=2.0cm, minimum height=0.72cm]
  (ROOT) at (2.70, 3.30) {\textbf{Root} $r$};
\node[font=\scriptsize\bfseries, text=clrRoot] at (2.70, 4.05)
  {stored on-chain};

% ── Edges ────────────────────────────────────────────────────
\draw[arr] (L0)--(H01); \draw[arr] (L1)--(H01);
\draw[arr] (L2)--(H23); \draw[arr] (L3)--(H23);
\draw[arr] (H01)--(ROOT); \draw[arr] (H23)--(ROOT);

% ── Verification path ─────────────────────────────────────────
\draw[prfarr, clrVerified!70!black]
  (L0.north) -- (H01.south west);
\draw[prfarr, clrComputed!80!black]
  (H01.north east) -- (ROOT.south west);

% ── Leaf formula ─────────────────────────────────────────────
\node[draw=gray!50, fill=gray!8, rounded corners=3pt,
      font=\scriptsize, align=center, minimum width=5.6cm,
      inner sep=4pt] at (2.70,-1.55)
  {$\ell_i = \mathrm{keccak256}(%
    \underbrace{\texttt{this}\|\texttt{chainid}}_{\text{context}} \|
    \underbrace{\texttt{issuer}\|\texttt{batchId}}_{\text{batch}} \|
    \underbrace{\texttt{docHash}_i}_{\text{diploma}})$};

\end{tikzpicture}%
\label{fig:merkle}%
}
\caption{UniVerify architecture and Merkle batch structure.}
\label{fig:overview}
\end{figure*}

% =============================================================
\section{Methodology}
\label{sec:method}

All experiments use Foundry~\cite{foundry2023}
(Solidity 0.8.20, optimiser at 200 runs).
Two measurement granularities are used throughout:
\emph{suite-level} gas (total per test, including setup)
and \emph{function-level} gas (Foundry's \texttt{--gas-report},
reporting min/avg/max over all invocations).
Seven versions of the per-diploma contract are studied
(Table~\ref{tab:versions}), each introducing exactly one change.
Two Merkle contracts are studied: \texttt{Baseline}
and \texttt{MerkleG} (four EVM-level hot-path optimisations).

\begin{table}[t]
\centering
\caption{Per-Diploma Contract Versions}
\label{tab:versions}
\setlength{\tabcolsep}{4pt}
\begin{tabular}{p{2.1cm} p{5.5cm}}
\toprule
Name & Key change \\
\midrule
\texttt{StringErrors}   & Baseline; \texttt{require(cond,"msg")}, full timestamp record \\
\texttt{TsRevoke}       & \texttt{revokedAt uint64} instead of \texttt{bool} — \textbf{regression} \\
\texttt{CustomErrors}   & Typed \texttt{error} declarations replace string messages \\
\texttt{ImmutableOwner} & \texttt{owner} as \texttt{immutable} (bytecode, not storage) \\
\texttt{Idempotent}     & \texttt{addIssuer}/\texttt{removeIssuer} silent on no-op \\
\texttt{MinimalRecord}  & Timestamps moved to event logs; record is \texttt{\{issuer, revoked\}} \\
\texttt{PackedStorage}  & Issuer + revoke bit packed into one \texttt{uint256} slot \\
\bottomrule
\end{tabular}
\end{table}

% =============================================================
\section{Gas Cost Results}

\subsection{Phase 1 — Per-Diploma Microoptimisations}
\label{sec:phase1}

Table~\ref{tab:phase1} reports function-level average gas.
Three effects are notable.

\emph{TsRevoke is a regression.}
A \texttt{uint64 revokedAt} timestamp occupies a new storage slot;
\texttt{revoke} cost jumps from 31,274 to 50,369 gas (+61\%).
Timestamps belong in event logs, not in storage.

\emph{Custom errors save bytecode.}
Replacing string messages with typed errors reduces deployment
from 464,211 to 409,497 gas (−11.9\%), because string literals
inflate bytecode while selectors are 4 bytes.

\emph{MinimalRecord achieves the best runtime.}
Removing all storage timestamps saves one cold \texttt{SSTORE}
per revocation.
Cumulative gains from \texttt{StringErrors} to \texttt{MinimalRecord}:
\begin{equation}
\Delta_\text{issue} \approx 14.1\%,\quad
\Delta_\text{revoke} \approx 7.8\%,\quad
\Delta_\text{deploy} \approx 19.6\%.
\label{eq:p1gains}
\end{equation}
Despite this, every version remains in the same cost class
($\sim$40–50k gas per diploma). The \emph{model} has not changed.

\begin{table}[t]
\centering
\caption{Phase 1 — Function-Level Gas (Avg)}
\label{tab:phase1}
\setlength{\tabcolsep}{4pt}
\begin{tabular}{lrrrr}
\toprule
Version & Deploy & \texttt{issue} & \texttt{revoke} \\
\midrule
\texttt{StringErrors}   & 464\,211 & 48\,458 & 31\,274 \\
\texttt{TsRevoke}       & 469\,181 & 50\,659 & \textbf{50\,369} \\
\texttt{CustomErrors}   & 409\,497 & 43\,367 & 38\,867 \\
\texttt{ImmutableOwner} & 403\,109 & 43\,370 & 38\,870 \\
\texttt{Idempotent}     & 418\,029 & 43\,377 & 38\,874 \\
\texttt{MinimalRecord}  & \textbf{373\,484} & \textbf{41\,607} & \textbf{28\,825} \\
\texttt{PackedStorage}  & 393\,596 & 43\,084 & 29\,302 \\
\bottomrule
\end{tabular}
\end{table}

\subsection{Phase 2 — Merkle Batch Architecture}
\label{sec:phase2}

A single \texttt{issueBatchRoot} call covers $N$ diplomas.
Per-diploma cost versus the per-diploma model:
\begin{equation}
\mathrm{gas}_\mathrm{Merkle}(N) = \frac{50{,}939}{N}
\quad \text{vs} \quad
\mathrm{gas}_\mathrm{legacy} = 48{,}200 = \text{const.}
\label{eq:amort}
\end{equation}
At $N = 100$: $509$ vs $48{,}200$ gas per diploma — a
\textbf{98.9\% reduction} (Fig.~\ref{fig:results}, right).
The Merkle contract costs $2.7\times$ more to deploy
(1,059,364 vs 393,596 gas), which breaks even at roughly $N \approx 22$
cumulative diplomas.

\subsection{Phase 3 — Merkle Hot-Path Optimisations}
\label{sec:phase3}

\texttt{MerkleG} applies four targeted changes to
\texttt{\_processProof()} and \texttt{\_merkleLeaf()}.
Table~\ref{tab:phase3} shows function-level results.
\texttt{issueBatchRoot} is unaffected (it does not enter either function).

\begin{table}[t]
\centering
\caption{Phase 3 — Baseline vs \texttt{MerkleG} (Function-Level Gas)}
\label{tab:phase3}
\setlength{\tabcolsep}{3.5pt}
\begin{tabular}{lrrrrrrc}
\toprule
& \multicolumn{3}{c}{\texttt{Baseline}} & \multicolumn{3}{c}{\texttt{MerkleG}} & $\Delta$Avg \\
\cmidrule(lr){2-4}\cmidrule(lr){5-7}
& Min & Avg & Max & Min & Avg & Max & \\
\midrule
Deploy         & \multicolumn{3}{c}{1{,}059{,}364}
               & \multicolumn{3}{c}{1{,}018{,}273} & $-3.9\%$ \\
\texttt{revokeFromBatch}
               & 54\,239 & 57\,229 & 61\,591
               & 54\,101 & 56\,443 & 59\,859 & $-1.4\%$ \\
\texttt{statusWithProof}
               & 10\,555 & 11\,890 & 13\,835
               & 10\,417 & 11\,104 & 12\,103 & \textbf{$-6.6\%$} \\
\texttt{merkleLeaf}
               & \multicolumn{3}{c}{866}
               & \multicolumn{3}{c}{824} & $-4.8\%$ \\
\bottomrule
\end{tabular}
\end{table}

\medskip
\noindent\textbf{Savings model.}
Running tests at proof depths $d \in \{0,1,4,8\}$
reveals a decomposable structure.
Let $N_\ell = 2^d$ be the number of leaves built in the test setup.
Total savings per test are:
\begin{equation}
\Delta(d) =
  \underbrace{N_\ell \cdot 42}_{\text{G4: per leaf}}
  +\;
  \underbrace{96}_{\text{G3: zero-checks}}
  +\;
  \underbrace{d \cdot 199}_{\text{G1+G2: per step}},
\label{eq:model}
\end{equation}
where 42~gas is saved per \texttt{merkleLeaf()} call (G4 assembly preimage),
96~gas is saved on the internal leaf call inside each function
(G3 skips redundant zero-guards), and 199~gas is saved per proof
step (G1 \texttt{unchecked} loop + G2 assembly keccak on EVM scratch space).
The model fits all four data points with zero residual error
(Fig.~\ref{fig:results}, left).

Note: the large absolute values at $d=8$ ($\approx$2.5M gas per test)
are dominated by the test harness building a 256-leaf tree in Solidity
— an artefact of simulation.
The actual on-chain call (\texttt{revokeFromBatch} with an 8-element
proof) costs $\approx$62k gas in both versions; the difference
is $\approx$1,732 gas.

% ── FIGURE 2: Phase 3 savings + amortisation (full-width) ───
\begin{figure*}[t]
\centering
\subfloat[G1--G4 savings model. Measured points (circles) match the
closed-form model~\eqref{eq:model} exactly (dashed line).]{%
\begin{tikzpicture}
\begin{axis}[
  width=8.3cm, height=5.0cm,
  xlabel={Proof depth $d$},
  ylabel={Gas saved (test-level)},
  xtick={0,1,4,8},
  ymode=log, log basis y=10,
  ymin=100, ymax=30000,
  grid=both,
  minor grid style={gray!12},
  major grid style={gray!35},
  legend style={at={(0.05,0.97)}, anchor=north west, font=\small},
  tick label style={font=\small},
  label style={font=\small},
  title style={font=\small},
  title={Phase 3: G1--G4 savings vs proof depth},
]
\addplot[mark=*, thick, color=blue!70!black, mark size=3pt]
  coordinates {(0,180)(1,421)(4,1606)(8,12484)};
\addlegendentry{Measured}
\addplot[mark=none, dashed, thick, color=red!65!black, domain=0:8, samples=50]
  {pow(2,x)*42 + 96 + x*199 + 42};
\addlegendentry{Model \eqref{eq:model}}
% Annotation at d=8
\node[font=\tiny, anchor=west, text=gray!70!black]
  at (axis cs:8.1, 12484) {12,484};
\end{axis}
\end{tikzpicture}%
\label{fig:p3savings}%
}\hfill
\subfloat[Per-diploma gas cost as a function of batch size $N$.
  Legacy cost is constant; Merkle cost decays as $1/N$.]{%
\begin{tikzpicture}
\begin{axis}[
  width=8.3cm, height=5.0cm,
  xlabel={Batch size $N$},
  ylabel={Gas per diploma},
  xmode=log, ymode=log,
  log basis x=10, log basis y=10,
  xmin=1, xmax=1000,
  ymin=30, ymax=200000,
  grid=both,
  minor grid style={gray!12},
  major grid style={gray!35},
  legend style={at={(0.97,0.97)}, anchor=north east, font=\small},
  tick label style={font=\small},
  label style={font=\small},
  title style={font=\small},
  title={Phase 2: amortisation of issuance cost},
]
\addplot+[mark=*, thick, color=blue!70!black, mark size=2.5pt]
  coordinates {(1,50939)(2,25470)(5,10188)(10,5094)(20,2547)
               (50,1019)(100,509)(200,255)(500,102)(1000,51)};
\addlegendentry{Merkle: $50939/N$}
\addplot+[mark=square*, thick, color=orange!75!black, dashed, mark size=2.5pt]
  coordinates {(1,48200)(2,48200)(5,48200)(10,48200)(20,48200)
               (50,48200)(100,48200)(200,48200)(500,48200)(1000,48200)};
\addlegendentry{Per-diploma (const.)}
% Annotations
\draw[dashed, gray!60] (axis cs:100,30)--(axis cs:100,509);
\node[font=\tiny, anchor=south, text=gray!70!black]
  at (axis cs:100,600) {$N=100$:\;$94.7\times$};
\draw[dashed, gray!60] (axis cs:1000,30)--(axis cs:1000,51);
\node[font=\tiny, anchor=south, text=gray!70!black]
  at (axis cs:1000,60) {$N=1000$:\;$945\times$};
\end{axis}
\end{tikzpicture}%
\label{fig:amort}%
}
\caption{Key quantitative results for Phase~3 (left) and Phase~2 (right).}
\label{fig:results}
\end{figure*}

% =============================================================
\section{Cost in USD}

Table~\ref{tab:usd} gives indicative USD costs at 40 gwei
and 1 ETH = \$\EthUsd.
At $N=100$, the full cost of certifying a graduating cohort
on-chain drops from $\approx$\$4.85 per diploma to $\approx$\$0.05 —
within practical reach for most universities.

\begin{table}[t]
\centering
\caption{Indicative Cost in USD (40 gwei, 1 ETH = \$\EthUsd)}
\label{tab:usd}
\setlength{\tabcolsep}{4pt}
\begin{tabular}{lrr}
\toprule
Operation & Gas & Cost (USD) \\
\midrule
\texttt{issue}, per diploma (legacy) & 48\,200 & \$\CostUsd{48200}{40} \\
\texttt{issueBatchRoot}, $N=1$ & 50\,939 & \$\CostUsd{50939}{40} \\
\texttt{issueBatchRoot}/diploma, $N=100$ & 509 & \$\CostUsd{509}{40} \\
\texttt{issueBatchRoot}/diploma, $N=1000$ & 51 & \$\CostUsd{51}{40} \\
\texttt{revokeFromBatch} (MerkleG, max) & 59\,859 & \$\CostUsd{59859}{40} \\
\texttt{statusWithProof} (view, free off-chain) & — & \$0.00 \\
Deployment (MerkleG) & 1\,018\,273 & \$\CostUsd{1018273}{40} \\
\bottomrule
\end{tabular}
\end{table}

% =============================================================
\section{Discussion and Conclusion}

\subsection{Hierarchy of Optimisation Impact}

The three phases reveal a clear hierarchy:
code-level microoptimisations deliver incremental gains (8–20\%);
the architectural shift delivers an amortisation gain of up to 99.9\%;
and EVM hot-path tuning provides a predictable, linear bonus.
Crucially, no amount of code-level polish in the per-diploma model
can bridge the gap to the Merkle batch model —
the two approaches belong to different complexity classes
in terms of per-unit cost.

\subsection{On-Chain vs Off-Chain Gas}

A recurring confusion in blockchain system measurements is the
conflation of \emph{simulation gas} with \emph{production gas}.
In our Phase 3 experiments, the test harness builds the Merkle tree
in Solidity (inside the EVM), which at depth 8 consumes
$\approx$2.2M gas — 87\% of the total test cost.
In production, this tree is built off-chain in ordinary code;
the on-chain cost for the same scenario is $\approx$113k gas.
Researchers reporting gas costs from test suites that simulate
off-chain work should always separate these two components.

\subsection{Conclusion}

We have shown that for a blockchain diploma registry,
the choice of data model dominates all other optimisation decisions.
The Merkle batch architecture reduces per-diploma issuance cost
by 98.9\% at $N=100$, compared to the 14\% improvement
achievable by exhaustive per-diploma code optimisation.
The four EVM hot-path optimisations (G1–G4) follow a closed-form
linear model and are recommended for any system with
sustained high-volume verification.
Future work will benchmark deeper trees ($d\geq12$),
calldata cost sensitivity, and alternative accumulator structures
such as sparse Merkle trees and Verkle trees.

% =============================================================
\begin{thebibliography}{9}

\bibitem{grech2017blockchain}
A.~Grech and A.~F.~Camilleri,
``Blockchain in Education,''
\emph{JRC Science for Policy Report}, European Commission, 2017.

\bibitem{alketbi2018blockchain}
A.~Alketbi, Q.~Nasir, and M.~A.~Talib,
``Blockchain for Government Services --- Use Cases, Security Benefits
and Challenges,''
in \emph{Proc.\ 15th Learn.\ Technol.\ Conf.}, 2018.

\bibitem{lizcano2020decentralised}
D.~Lizcano, J.~A.~Lara, B.~White, and S.~Aljawarneh,
``Decentralised technology for educational credentials,''
\emph{J.\ Comput.\ High.\ Educ.}, vol.~32, pp.~9--49, 2020.

\bibitem{merkle1987}
R.~C.~Merkle,
``A Digital Signature Based on a Conventional Encryption Function,''
in \emph{Advances in Cryptology --- CRYPTO '87},
LNCS, vol.~293, pp.~369--378, 1988.

\bibitem{eip2929}
V.~Buterin et~al.,
``EIP-2929: Gas Cost Increases for State Access Opcodes,''
\emph{Ethereum Improvement Proposals}, 2020.

\bibitem{foundry2023}
Paradigm,
``Foundry: Ethereum Application Development Toolkit,''
\url{https://github.com/foundry-rs/foundry}, 2023.

\bibitem{yellowpaper}
G.~Wood,
``Ethereum: A Secure Decentralised Generalised Transaction Ledger,''
\emph{Ethereum Yellow Paper}, rev.\ 2023.

\end{thebibliography}

\end{document}
